% Introduction and Related Work draft

\section{Introduction}

Normal aging is accompanied by declines in higher cognitive abilities that are disproportionately pronounced for functions mediated by the dorsolateral prefrontal cortex (dlPFC), including spatial working memory. In the rhesus monkey, which provides the closest primate model of human cortical aging, behavioral testing consistently reveals working memory deficits in aged animals even in the absence of overt neurodegeneration \citep{chang2005firing,luebke2010area46}. Unlike classical neurodegenerative conditions, this decline emerges without frank neuron loss, pointing instead to subtler, distributed changes in neuronal structure and function as the substrates of impaired cognition.

At the cellular level, layer 3 pyramidal neurons of area 46 undergo a distinctive combination of sublethal changes with age. Dendritic arbors regress, spine and synapse counts fall by roughly thirty percent in layers 2/3, and intrinsic excitability rises, producing a characteristic increase in action potential firing rates \citep{peters2008synapses,chang2005firing,coskren2015aging,dickstein2013spines}. These changes have been linked to cognitive decline both empirically and through computational models that explore how altered f--I curves and synaptic densities modify persistent activity in working memory networks \citep{ibanez2020network}. A parallel line of ultrastructural research, however, has revealed an equally striking set of age-related changes affecting the myelin sheaths that insulate the axons of these same pyramidal neurons. In area 46 of aged monkeys, 3--6\% of myelin sheaths display dystrophic alterations including splitting of the major dense line, ``balloons'' of separated intraperiod lines, and redundant myelin, and the frequency of these alterations correlates significantly with composite cognitive impairment across subjects \citep{peters2002aging}. Transverse sections reveal a 90\% increase in the number of paranodal profiles in aged versus young monkeys, consistent with the addition of shorter internodes through remyelination, and similar changes are documented in the fornix, cingulum and corpus callosum \citep{peters2003remyelination,peters2010fornix,bowley2010cingulate}. More recent cellular studies indicate that aging compromises the maturation of oligodendrocyte precursor cells and the capacity for efficient remyelination, further implicating myelin as a mechanistic locus of age-related dysfunction \citep{dimovasili2022oligodendrocyte,go2020vesicles}.

Despite this rich descriptive base, the functional consequences of age-related myelin dystrophy for prefrontal computation remain poorly quantified. Classical biophysical theory holds that myelin thickness, internode length, and axon diameter jointly determine conduction velocity \citep{waxman1980determinants}, and modeling work has explored how internodal tight junctions and the precise geometry of nodes of Ranvier modulate the safety factor of propagation \citep{devaux2008tight,smith2024nodal}. Existing models of demyelination and remyelination, however, have typically considered either condition in isolation, and have largely been calibrated to peripheral axons or to multiple sclerosis rather than to primate cortical aging \citep{stephanova2003extracellular,coggan2015demyelination,naud2019demyelination}. Computational work on remyelinated segments indicates that shorter and thinner new internodes can slow conduction even when all previously demyelinated segments are covered, an effect that is amplified by transitions between normal and short internodes \citep{lasiene2008partial,powers2012remyelination,scurfield2018internodal}. What is missing is a pipeline that (i) embeds biophysically detailed axons in realistic, empirically constrained pyramidal neurons, (ii) systematically varies the proportion of segments affected and the fraction of lamellae removed or restored, and (iii) propagates the resulting heterogeneous action potential failures into a cortical network capable of sustaining working memory.

The present study addresses this gap by coupling two models at different scales. At the single-neuron scale, we construct a cohort of fifty multicompartment dlPFC layer 3 pyramidal neurons with detailed axons, built by Latin Hypercube Sampling (LHS) over eight axon parameters and constrained to match empirically observed somatic firing rates, conduction velocities, and g-ratios \citep{rumbell2016automated,coskren2015aging}. Each axon includes 101 nodes separating 100 myelinated segments, each composed of paranodes, juxtaparanodes, an internode, and tight junctions, with electrical properties drawn from prior myelinated-axon models \citep{devaux2008tight,scurfield2018internodal}. Demyelination is implemented by removing a specified percentage of lamellae from a specified subset of segments, and remyelination by replacing affected segments with two shorter, thinner segments separated by a new node. At the network scale, we adapt a recently proposed spiking network with facilitating recurrent excitation \citep{hansel2013short} to simulate an oculomotor delayed response task based on the classical paradigm of \citet{funahashi1989mnemonic}, and we transfer action potential failure distributions measured in the single-neuron cohort into the network as probabilistic transmission failures. This makes it possible to ask, for the first time, whether myelin alterations at empirically observed levels are sufficient to produce quantitative deficits in memory duration and precision.

This work makes four contributions. First, we show that complete demyelination of even 25\% of segments slows axonal conduction by $38 \pm 10\%$ and causes $35 \pm 13\%$ of action potentials to fail, with further increases as more segments are affected. Second, we demonstrate that remyelination with sufficient new lamellae can largely rescue transmission, but that partial remyelination leaves substantial failure rates. Third, we map these single-neuron effects onto working memory performance in a 20{,}000-neuron spiking network, showing that propagation delays alone have negligible effect while AP failures reduce memory duration and increase diffusion in a graded, empirically matched fashion. Fourth, by quantifying memory performance as a function of the percentage of normal or new myelin sheaths in heterogeneous axon populations, we link our simulations directly to the ultrastructural ranges reported by \citet{peters2002aging} and \citet{peters2003remyelination}, providing a mechanistic account of how observed levels of myelin dystrophy could generate the working memory decline that characterizes cognitive aging in the primate.

\section{Related Work}

\paragraph{Age-related changes in prefrontal neurons and circuits.}
Pyramidal neurons in aged rhesus monkey dlPFC show multiple sublethal alterations that have been linked to working memory decline. Patch-clamp recordings from layer 2/3 pyramidal cells demonstrate elevated intrinsic excitability and action potential firing rates, with a U-shaped relationship between firing rate and cognitive performance \citep{chang2005firing}. Ultrastructural quantification reveals a loss of approximately 30\% of asymmetric and symmetric synapses from layers 2/3 with age, and this loss correlates with the cognitive impairment index of the same animals \citep{peters2008synapses,dickstein2013spines}. Morphometric reconstructions combined with biophysical modeling indicate that dendritic arbor regression together with altered passive and active membrane properties can reproduce the observed physiological phenotype \citep{coskren2015aging,rumbell2016automated}. Building on these empirical findings, \citet{ibanez2020network} showed in a bump-attractor network that increased firing rates and synapse loss can together impair persistent activity during an oculomotor delayed response task; crucially, that work did not incorporate any changes to myelin structure, leaving open the functional role of the myelin dystrophies documented in the same brain regions. Our study fills this specific gap by introducing a mechanistic link from sheath-level alterations to the population dynamics that encode spatial working memory.

\paragraph{Myelin alterations in the aging primate brain.}
Electron microscopic studies across multiple cortical and subcortical tracts have established that normal aging in the rhesus monkey produces a distinctive constellation of myelin alterations. In area 46, roughly 3--6\% of sheaths exhibit splitting of the major dense line, balloons, or redundant myelin, and the frequency of these features correlates with cognitive impairment both across and within aged monkeys \citep{peters2002aging}. In aged area 46 specifically, the number of paranodal profiles is 90\% higher than in young animals, indicating more and shorter internodes and consistent with ongoing remyelination across the lifespan \citep{peters2003remyelination}. Comparable patterns appear in the corpus callosum, cingulate bundle, and fornix, with fiber tract--specific correlations to cognitive domains \citep{bowley2010cingulate,peters2010fornix}. The cellular basis of these changes is increasingly clear: aging reduces the numbers of newly differentiating myelinating oligodendrocytes and impairs the transcriptional programme underlying oligodendrocyte precursor cell maturation, providing a mechanism for inefficient remyelination \citep{dimovasili2022oligodendrocyte,go2020vesicles}.

\paragraph{Computational models of myelinated axon function.}
Classical cable-theoretic analysis shows that for myelinated fibers, conduction velocity depends jointly on axon diameter, myelin thickness, and internode length, with an optimal g-ratio for maximal propagation speed \citep{waxman1980determinants}. More recent models have incorporated biophysically detailed descriptions of nodes, paranodes, juxtaparanodes, internodes, and myelin tight junctions, and have shown that tight junctions contribute substantially to internodal insulation in small CNS fibers \citep{devaux2008tight}. Developmental work links the final geometry of nodes of Ranvier to conduction speed and timing precision \citep{smith2024nodal}. Models of demyelination typically remove lamellae from a single fiber or segment to examine effects on the compound action potential \citep{stephanova2003extracellular,coggan2015demyelination,naud2019demyelination}. Remyelination has been modeled by introducing shorter or thinner segments, with the clear prediction that transitions between normal and remyelinated internodes depress conduction beyond what mean internode length alone would predict \citep{scurfield2018internodal,powers2012remyelination,lasiene2008partial}. Our work combines these disparate modeling ingredients within a single biologically constrained cohort and, for the first time, couples them to a network-level memory task.

\paragraph{Network models of spatial working memory.}
Since the seminal recordings of \citet{funahashi1989mnemonic}, spatial working memory has been conceptualized as the maintenance of persistent, spatially tuned activity in a ring of prefrontal excitatory neurons. The canonical attractor model of \citet{compte2000synaptic} showed that strong NMDA-mediated recurrent excitation, balanced by lateral inhibition, yields a stable activity ``bump'' whose diffusion captures the behavioral precision of delayed-response performance. \citet{hansel2013short} extended this framework by showing that short-term synaptic facilitation in recurrent excitation accounts for the irregular firing observed during persistent activity and supports robust memory storage; we adopt this facilitating recurrent architecture as our network substrate. Bump-attractor dynamics have since been empirically validated against prefrontal single-unit activity and behavioral errors \citep{wimmer2014bump,panichello2019error}, and extended to account for serial biases via interacting persistent and activity-silent representations \citep{barbosa2020interplay}. An alternative class of activity-silent models invokes short-term synaptic plasticity as the memory trace \citep{mongillo2008synaptic,wolff2015hidden}; we use this class to test the generality of our conclusions. Our contribution is to couple these well-established network dynamics to biophysically grounded axonal pathology, producing a direct mapping from measurable ultrastructural statistics to measurable behavioral indices of memory duration and precision.
